% answers.tex —— 2014 年·安徽卷·文科 · 参考答案与详细解析片段
% 由 tools/generate_projects.py 依来源总账自 2014/2014_数学真题_安徽文_试题.md + 2014/2014_数学真题_安徽文_解析正文.md + 2014/2014_数学真题_安徽文_答案要点.md 生成，勿手改（改源 MD 后重跑）


\begin{answer}
\gksolhead{1}{}

\textbf{题目}：

设 $\mathrm i$ 是虚数单位，复数 $\mathrm i^3+\dfrac{2\mathrm i}{1+\mathrm i}=$\mbox{（　　）}

（A）$-\mathrm i$；（B）$\mathrm i$；（C）$-1$；（D）$1$。

\Needspace{4\baselineskip}
\textbf{解析}

\Needspace{5\baselineskip}
由 $\mathrm i^2=-1$，先分别化简两项：
\nopagebreak[4]
\[\mathrm i^3=\mathrm i^2\cdot\mathrm i=-\mathrm i,\qquad
\frac{2\mathrm i}{1+\mathrm i}=\frac{2\mathrm i(1-\mathrm i)}{(1+\mathrm i)(1-\mathrm i)}=\frac{2\mathrm i+2}{2}=1+\mathrm i .\]
\Needspace{5\baselineskip}
相加得
\nopagebreak[4]
\[\mathrm i^3+\frac{2\mathrm i}{1+\mathrm i}=-\mathrm i+(1+\mathrm i)=1 .\]

选 D。

\gkanswer{答案：D。$\mathrm i^3=-\mathrm i$，$\dfrac{2\mathrm i}{1+\mathrm i}=\mathrm i(1-\mathrm i)=1+\mathrm i$，相加得 $1$。}

\end{answer}
\begin{answer}
\gksolhead{2}{}

\textbf{题目}：

命题“$\forall x\in\mathbf R$，$|x|+x^2\geqslant0$”的否定是\mbox{（　　）}

（A）$\forall x\in\mathbf R$，$|x|+x^2<0$；（B）$\forall x\in\mathbf R$，$|x|+x^2\leqslant0$；（C）$\exists x_0\in\mathbf R$，$|x_0|+x_0^2<0$；（D）$\exists x_0\in\mathbf R$，$|x_0|+x_0^2\geqslant0$。

\Needspace{4\baselineskip}
\textbf{解析}

\Needspace{5\baselineskip}
原命题是全称命题“$\forall x\in\mathbf R$，$p(x)$”，其中 $p(x)$ 为 $|x|+x^2\geqslant0$。全称命题的否定是特称命题，且结论要取反：
\nopagebreak[4]
\[\neg\bigl(\forall x\in\mathbf R,\ |x|+x^2\geqslant0\bigr)
\iff \exists x_0\in\mathbf R,\ |x_0|+x_0^2<0 .\]
（A）（B）仍用全称量词；（D）没有把“$\geqslant$”改成“$<$”。

选 C。

\gkanswer{答案：C。全称命题的否定是特称命题，同时把结论“$\geqslant0$”改为“$<0$”。}

\end{answer}
\begin{answer}
\gksolhead{3}{}

\textbf{题目}：

抛物线 $y=\dfrac14x^2$ 的准线方程是\mbox{（　　）}

（A）$y=-1$；（B）$y=-2$；（C）$x=-1$；（D）$x=-2$。

\Needspace{4\baselineskip}
\textbf{解析}

\Needspace{5\baselineskip}
把 $y=\dfrac14x^2$ 写成标准形式
\nopagebreak[4]
\[x^2=4y .\]
\Needspace{5\baselineskip}
与 $x^2=2py$ 对比得 $2p=4$，即 $p=2$。开口向上、顶点在原点的抛物线，其准线为
\nopagebreak[4]
\[y=-\frac p2=-1 .\]
焦点为 $(0,1)$，准线是水平线，故（C）（D）不成立。

选 A。

\gkanswer{答案：A。抛物线化为 $x^2=4y$，$2p=4$，故准线为 $y=-1$。}

\end{answer}
\begin{answer}
\gksolhead{4}{}

\textbf{题目}：

如图所示，程序框图（算法流程图）的输出结果是\mbox{（　　）}

（A）$34$；（B）$55$；（C）$78$；（D）$89$。
\nopagebreak[4]
\begin{center}% ============================================================
% f2014_anhui_arts_q04_flowchart.tex
% 2014 安徽卷（文科）第 4 题：程序框图
% 来源 PDF：普通高考/2014/2014安徽文.pdf
% 原页：第 1 页（A3 横排）左栏，印面 859，第 4 题题干下方
% 结构：开始 → 赋初值 x=1、y=1 → z=x+y → 判断 z≤50；
%       是：x=y → y=z → 返回 z=x+y 之前；否：输出 z → 结束。
% ============================================================
\begin{tikzpicture}[>=Stealth, line width=0.7pt, font=\small,
  x=1cm, y=1cm, line cap=round, line join=round]

  % 开始
  \draw[rounded corners=6pt] (-0.95,0.42) rectangle (0.95,-0.42);
  \node at (0,0) {开始};
  \draw[->] (0,-0.42) -- (0,-0.85);

  % 赋初值
  \draw (-1.55,-0.85) rectangle (1.55,-1.70);
  \node at (0,-1.28) {$x=1$，$y=1$};
  \draw[->] (0,-1.70) -- (0,-2.10);

  % z=x+y
  \draw (-1.30,-2.10) rectangle (1.30,-2.95);
  \node at (0,-2.52) {$z=x+y$};
  \draw[->] (0,-2.95) -- (0,-3.40);

  % 判断框
  \draw (0,-3.40) -- (1.75,-4.25) -- (0,-5.10) -- (-1.75,-4.25) -- cycle;
  \node at (0,-4.25) {$z\leqslant 50$};

  % 否：向右输出
  \draw[->] (1.75,-4.25) -- (3.30,-4.25);
  \node at (2.20,-3.95) {否};
  \draw (3.30,-3.82) rectangle (5.50,-4.68);
  \node at (4.40,-4.25) {输出 $z$};
  \draw[->] (4.40,-4.68) -- (4.40,-5.05);

  % 结束
  \draw[rounded corners=6pt] (3.45,-5.05) rectangle (5.35,-5.90);
  \node at (4.40,-5.47) {结束};

  % 是：向左进入循环体
  \draw[->] (-1.75,-4.25) -- (-3.30,-4.25) -- (-3.30,-4.75);
  \node at (-2.30,-3.95) {是};
  \draw (-4.45,-4.75) rectangle (-2.15,-5.60);
  \node at (-3.30,-5.17) {$x=y$};

  \draw[->] (-3.30,-5.60) -- (-3.30,-5.95);
  \draw (-4.45,-5.95) rectangle (-2.15,-6.80);
  \node at (-3.30,-6.37) {$y=z$};

  % 回边：向下、向左、向上，再水平汇入 z=x+y 与判断框之间
  \draw (-3.30,-6.80) -- (-3.30,-7.25) -- (-6.30,-7.25) -- (-6.30,-3.18) -- (0,-3.18);
  \draw[->] (0,-3.18) -- (0,-3.40);

\end{tikzpicture}
\end{center}
\nopagebreak[4]
\Needspace{4\baselineskip}
\textbf{解析}

按框图顺序追踪：初值 $x=1$，$y=1$；每轮先算 $z=x+y$，再判断 $z\leqslant50$，成立则令 $x$ 取旧的 $y$、$y$ 取 $z$，然后返回。

\begin{center}
\begin{tabular}{llllll}
\toprule
轮次 & 判断前 $x$ & 判断前 $y$ & $z=x+y$ & $z\leqslant50$ & 新 $(x,y)$ \\
\midrule
1 & $1$ & $1$ & $2$ & 是 & $(1,2)$ \\
2 & $1$ & $2$ & $3$ & 是 & $(2,3)$ \\
3 & $2$ & $3$ & $5$ & 是 & $(3,5)$ \\
4 & $3$ & $5$ & $8$ & 是 & $(5,8)$ \\
5 & $5$ & $8$ & $13$ & 是 & $(8,13)$ \\
6 & $8$ & $13$ & $21$ & 是 & $(13,21)$ \\
7 & $13$ & $21$ & $34$ & 是 & $(21,34)$ \\
8 & $21$ & $34$ & $55$ & 否 & 输出 $55$ \\
\bottomrule
\end{tabular}
\end{center}

即 $z$ 依次为斐波那契型数列 $2,3,5,8,13,21,34,55$，第一个大于 $50$ 的是 $55$，故输出 $55$。

选 B。

\gkanswer{答案：B。$z$ 依次取 $2,3,5,8,13,21,34,55$，$55>50$ 时输出。}

\end{answer}
\begin{answer}
\gksolhead{5}{}

\textbf{题目}：

设 $a=\log_37$，$b=2^{1.1}$，$c=0.8^{3.1}$，则\mbox{（　　）}

（A）$b<a<c$；（B）$c<a<b$；（C）$c<b<a$；（D）$a<c<b$。

\Needspace{4\baselineskip}
\textbf{解析}

\Needspace{5\baselineskip}
分别估计三个数与 $0$，$1$，$2$ 的大小：
\nopagebreak[4]
\[a=\log_37,\qquad 3^1<7<3^2\ \Longrightarrow\ 1<a<2 ;\]
\nopagebreak[4]
\[b=2^{1.1}>2^1=2 ;\]
\nopagebreak[4]
\[c=0.8^{3.1}<0.8^0=1,\qquad c>0 .\]
于是 $c<1<a<2<b$，即 $c<a<b$。

选 B。

\gkanswer{答案：B。$a=\log_37\in(1,2)$，$b=2^{1.1}>2$，$c=0.8^{3.1}\in(0,1)$，故 $c<a<b$。}

\end{answer}
\begin{answer}
\gksolhead{6}{}

\textbf{题目}：

过点 $P(-\sqrt3,-1)$ 的直线 $l$ 与圆 $x^2+y^2=1$ 有公共点，则直线 $l$ 的倾斜角的取值范围是\mbox{（　　）}

（A）$\left(0,\dfrac\pi6\right]$；（B）$\left(0,\dfrac\pi3\right]$；（C）$\left[0,\dfrac\pi6\right]$；（D）$\left[0,\dfrac\pi3\right]$。

\Needspace{4\baselineskip}
\textbf{解析}

\Needspace{5\baselineskip}
点 $P(-\sqrt3,-1)$ 在圆 $x^2+y^2=1$ 外（$3+1=4>1$）。设 $l$ 的斜率为 $k$，则
\nopagebreak[4]
\[l:\ y+1=k(x+\sqrt3)\ \Longrightarrow\ kx-y+\sqrt3k-1=0 .\]
\Needspace{5\baselineskip}
$l$ 与圆有公共点，等价于圆心 $(0,0)$ 到 $l$ 的距离不超过半径 $1$：
\nopagebreak[4]
\[\frac{|\sqrt3k-1|}{\sqrt{k^2+1}}\leqslant1
\ \Longrightarrow\ 3k^2-2\sqrt3k+1\leqslant k^2+1
\ \Longrightarrow\ 2k^2-2\sqrt3k\leqslant0,\]
\Needspace{5\baselineskip}
即
\nopagebreak[4]
\[k\left(k-\sqrt3\right)\leqslant0\ \Longrightarrow\ 0\leqslant k\leqslant\sqrt3 .\]
当 $k$ 从 $0$ 增到 $\sqrt3$ 时，倾斜角 $\alpha=\arctan k$ 从 $0$ 增到 $\dfrac\pi3$。再看斜率不存在的竖直直线 $x=-\sqrt3$：它到圆心的距离为 $\sqrt3>1$，与圆无公共点，不产生补充的角。

故倾斜角取值范围是 $\left[0,\dfrac\pi3\right]$。

选 D。

\gkanswer{答案：D。设斜率为 $k$，圆心到 $l$ 的距离 $\leqslant1$ 给出 $0\leqslant k\leqslant\sqrt3$，倾斜角属 $\left[0,\dfrac\pi3\right]$。}

\end{answer}
\begin{answer}
\gksolhead{7}{}

\textbf{题目}：

若将函数 $f(x)=\sin2x+\cos2x$ 的图象向右平移 $\varphi$ 个单位，所得图象关于 $y$ 轴对称，则 $\varphi$ 的最小正值是\mbox{（　　）}

（A）$\dfrac\pi8$；（B）$\dfrac\pi4$；（C）$\dfrac{3\pi}8$；（D）$\dfrac{3\pi}4$。

\Needspace{4\baselineskip}
\textbf{解析}

\Needspace{5\baselineskip}
先合并：
\nopagebreak[4]
\[f(x)=\sin2x+\cos2x=\sqrt2\sin\left(2x+\frac\pi4\right).\]
\Needspace{5\baselineskip}
向右平移 $\varphi$（$\varphi>0$）个单位，得
\nopagebreak[4]
\[g(x)=f(x-\varphi)=\sqrt2\sin\left(2x-2\varphi+\frac\pi4\right).\]
\Needspace{5\baselineskip}
$g(x)$ 的图象关于 $y$ 轴对称，等价于 $g$ 为偶函数。正弦型函数为偶函数时其相位必须是 $\dfrac\pi2+k\pi$：
\nopagebreak[4]
\[-2\varphi+\frac\pi4=\frac\pi2+k\pi\ \Longrightarrow\ \varphi=-\frac\pi8-\frac{k\pi}2\quad(k\in\mathbf Z).\]
取 $k=-1$ 得 $\varphi=\dfrac{3\pi}8>0$；$k=0$ 时 $\varphi=-\dfrac\pi8<0$，$k\leqslant-2$ 时 $\varphi\geqslant\dfrac{7\pi}8$ 更大。故最小正值为 $\dfrac{3\pi}8$。

验证：$\varphi=\dfrac{3\pi}8$ 时 $g(x)=\sqrt2\sin\left(2x-\dfrac\pi2\right)=-\sqrt2\cos2x$，确为偶函数。

选 C。

\gkanswer{答案：C。$f(x)=\sqrt2\sin\left(2x+\dfrac\pi4\right)$，平移后相位为 $\dfrac\pi4-2\varphi$，关于 $y$ 轴对称需 $\dfrac\pi4-2\varphi=\dfrac\pi2+k\pi$，最小正值 $\varphi=\dfrac{3\pi}8$。}

\end{answer}
\begin{answer}
\gksolhead{8}{}

\textbf{题目}：

一个多面体的三视图如图所示，则该多面体的体积为\mbox{（　　）}

（A）$\dfrac{23}3$；（B）$\dfrac{47}6$；（C）$6$；（D）$7$。
\nopagebreak[4]
\begin{center}% ============================================================
% f2014_anhui_arts_q08_three_views.tex
% 2014 安徽卷（文科）第 8 题：三视图
% 来源 PDF：普通高考/2014/2014安徽文.pdf
% 原页：第 1 页（A3 横排）中栏上部，印面 859，第 8 题题干上方
% 说明：三个视图都是边长 2 的正方形（边标注 1 与 1）。
%       正视图：实线 (0,1)-(1,0)，虚线 (1,2)-(2,1)；
%       侧视图：虚线 (0,1)-(1,2)，实线 (1,0)-(2,1)；
%       俯视图：实线 (1,2)-(2,1)，虚线 (0,1)-(1,0)。
% ============================================================
\begin{tikzpicture}[>=Stealth, line width=0.7pt, font=\small,
  x=1.55cm, y=1.55cm, line cap=round, line join=round,
  ext/.style={line width=0.4pt},
  dim/.style={<->, line width=0.45pt}]

  % ---------- 正（主）视图 ----------
  \draw[ext] (0,0) -- (2,0) -- (2,2) -- (0,2) -- cycle;
  \draw (0,1) -- (1,0);
  \draw[dashed] (1,2) -- (2,1);
  % 上方 1 + 1
  \draw[ext] (0,2) -- (0,2.55);
  \draw[ext] (1,2) -- (1,2.30);
  \draw[ext] (2,2) -- (2,2.55);
  \draw[dim] (0,2.42) -- (0.92,2.42) node[midway,above,inner sep=1pt] {$1$};
  \draw[dim] (1.08,2.42) -- (2,2.42) node[midway,above,inner sep=1pt] {$1$};
  % 左侧 1 + 1
  \draw[ext] (0,2) -- (-0.55,2);
  \draw[ext] (0,1) -- (-0.30,1);
  \draw[ext] (0,0) -- (-0.55,0);
  \draw[dim] (-0.42,1.08) -- (-0.42,2) node[midway,left,inner sep=1pt] {$1$};
  \draw[dim] (-0.42,0) -- (-0.42,0.92) node[midway,left,inner sep=1pt] {$1$};
  \node at (1,-0.62) {正（主）视图};

  % ---------- 侧（左）视图 ----------
  \begin{scope}[xshift=4.75cm]
    \draw[ext] (0,0) -- (2,0) -- (2,2) -- (0,2) -- cycle;
    \draw[dashed] (0,1) -- (1,2);
    \draw (1,0) -- (2,1);
    % 上方 1 + 1
    \draw[ext] (0,2) -- (0,2.55);
    \draw[ext] (1,2) -- (1,2.30);
    \draw[ext] (2,2) -- (2,2.55);
    \draw[dim] (0,2.42) -- (0.92,2.42) node[midway,above,inner sep=1pt] {$1$};
    \draw[dim] (1.08,2.42) -- (2,2.42) node[midway,above,inner sep=1pt] {$1$};
    % 左侧 1 + 1
    \draw[ext] (0,2) -- (-0.55,2);
    \draw[ext] (0,1) -- (-0.30,1);
    \draw[ext] (0,0) -- (-0.55,0);
    \draw[dim] (-0.42,1.08) -- (-0.42,2) node[midway,left,inner sep=1pt] {$1$};
    \draw[dim] (-0.42,0) -- (-0.42,0.92) node[midway,left,inner sep=1pt] {$1$};
    \node at (1,-0.62) {侧（左）视图};
  \end{scope}

  % ---------- 俯视图 ----------
  \begin{scope}[yshift=-5.15cm]
    \draw[ext] (0,0) -- (2,0) -- (2,2) -- (0,2) -- cycle;
    \draw (1,2) -- (2,1);
    \draw[dashed] (0,1) -- (1,0);
    % 下方 1 + 1
    \draw[ext] (0,0) -- (0,-0.55);
    \draw[ext] (1,0) -- (1,-0.30);
    \draw[ext] (2,0) -- (2,-0.55);
    \draw[dim] (0,-0.42) -- (0.92,-0.42) node[midway,below,inner sep=1pt] {$1$};
    \draw[dim] (1.08,-0.42) -- (2,-0.42) node[midway,below,inner sep=1pt] {$1$};
    % 左侧 1 + 1
    \draw[ext] (0,2) -- (-0.55,2);
    \draw[ext] (0,1) -- (-0.30,1);
    \draw[ext] (0,0) -- (-0.55,0);
    \draw[dim] (-0.42,1.08) -- (-0.42,2) node[midway,left,inner sep=1pt] {$1$};
    \draw[dim] (-0.42,0) -- (-0.42,0.92) node[midway,left,inner sep=1pt] {$1$};
    \node at (1,-0.98) {俯视图};
  \end{scope}

\end{tikzpicture}
\end{center}
\nopagebreak[4]
\Needspace{4\baselineskip}
\textbf{解析}

三个视图的外轮廓都是边长 $2$ 的正方形，说明多面体可由棱长 $2$ 的正方体切去若干角块得到；每个视图内部各有一条实线与一条虚线，说明切去的是正方体上两个相对顶点处的三棱锥。

取正方体 $[0,2]^3$，切去顶点 $O(0,0,0)$ 处的三棱锥，截面过 $(1,0,0)$，$(0,1,0)$，$(0,0,1)$；再切去顶点 $(2,2,2)$ 处的三棱锥，截面过 $(1,2,2)$，$(2,1,2)$，$(2,2,1)$。逐个视图核对：

\begin{itemize}
\item 正视图（沿 $y$ 方向投影）：$O$ 处截面投影为 $(0,1)$ 到 $(1,0)$ 的线段，朝前可见，画实线；$(2,2,2)$ 处截面投影为 $(1,2)$ 到 $(2,1)$ 的线段，被遮挡，画虚线；
\item 侧视图（沿 $x$ 方向投影，水平轴取反向）：对应的两条投影线恰为 $(1,0)$ 到 $(2,1)$ 的实线与 $(1,2)$ 到 $(0,1)$ 的虚线；
\item 俯视图（沿 $z$ 方向投影）：$O$ 处截面在底面，被遮挡，投影为 $(0,1)$ 到 $(1,0)$ 的虚线；$(2,2,2)$ 处截面朝上可见，投影为 $(1,2)$ 到 $(2,1)$ 的实线。
\end{itemize}

\Needspace{5\baselineskip}
与题给三视图一一吻合，故多面体就是正方体挖去这两个三棱锥。每个三棱锥的三条侧棱两两垂直、长均为 $1$，体积
\nopagebreak[4]
\[V_1=\frac16\times1\times1\times1=\frac16 .\]
\Needspace{5\baselineskip}
于是
\nopagebreak[4]
\[V=2^3-2\times\frac16=8-\frac13=\frac{23}3 .\]

选 A。

\gkanswer{答案：A。该多面体是棱长 $2$ 的正方体挖去两个三侧棱两两垂直且长均为 $1$ 的四面体（正方体一角的角块），$8-2\times\dfrac16=\dfrac{23}3$。}

\end{answer}
\begin{answer}
\gksolhead{9}{}

\textbf{题目}：

若函数 $f(x)=|x+1|+|2x+a|$ 的最小值为 $3$，则实数 $a$ 的值为\mbox{（　　）}

（A）$5$ 或 $8$；（B）$-1$ 或 $5$；（C）$-1$ 或 $-4$；（D）$-4$ 或 $8$。

\Needspace{4\baselineskip}
\textbf{解析}

\Needspace{5\baselineskip}
记 $f(x)=|x+1|+2\left|x+\dfrac a2\right|$，它是分段线性函数，折点只有 $x=-1$ 与 $x=-\dfrac a2$，最小值必在折点处取得：
\nopagebreak[4]
\[f(-1)=\left|a-2\right|,\qquad
f\left(-\frac a2\right)=\left|1-\frac a2\right| .\]
令最小值为 $3$：

\begin{itemize}
\item 由 $|a-2|=3$ 得 $a=5$ 或 $a=-1$；
\item 由 $\left|1-\dfrac a2\right|=3$ 得 $a=8$ 或 $a=-4$。
\end{itemize}

代回检验哪个折点真正给出最小值：

\begin{itemize}
\item $a=5$：折点为 $-1$ 与 $-\dfrac52$，$f\left(-\dfrac52\right)=\left|-\dfrac32\right|=\dfrac32$，$f(-1)=3$，最小值是 $\dfrac32\neq3$，舍去；
\item $a=-1$：折点为 $-1$ 与 $\dfrac12$，$f\left(\dfrac12\right)=\left|1+\dfrac12\right|=\dfrac32$，最小值不是 $3$，舍去；
\item $a=-4$：折点为 $-1$ 与 $2$，$f(2)=|3|+0=3$，且 $x>2$ 时 $f$ 递增、$-1<x<2$ 时 $f$ 递减，最小值确为 $3$；
\item $a=8$：折点为 $-1$ 与 $-4$，$f(-4)=0+|-3|=3$，同理最小值确为 $3$。
\end{itemize}

故 $a=-4$ 或 $a=8$。

选 D。

\gkanswer{答案：D。最小值在折点取得，$|a-2|=3$ 或 $\left|1-\dfrac a2\right|=3$，逐一代回检验得 $a=-4$ 或 $a=8$。}

\end{answer}
\begin{answer}
\gksolhead{10}{}

\textbf{题目}：

设 $\boldsymbol a$，$\boldsymbol b$ 为非零向量，$|\boldsymbol b|=2|\boldsymbol a|$，两组向量 $\boldsymbol x_1$，$\boldsymbol x_2$，$\boldsymbol x_3$，$\boldsymbol x_4$ 和 $\boldsymbol y_1$，$\boldsymbol y_2$，$\boldsymbol y_3$，$\boldsymbol y_4$ 均由 $2$ 个 $\boldsymbol a$ 和 $2$ 个 $\boldsymbol b$ 排列而成，若 $\boldsymbol x_1\cdot\boldsymbol y_1+\boldsymbol x_2\cdot\boldsymbol y_2+\boldsymbol x_3\cdot\boldsymbol y_3+\boldsymbol x_4\cdot\boldsymbol y_4$ 所有可能取值中的最小值为 $4|\boldsymbol a|^2$，则 $\boldsymbol a$ 与 $\boldsymbol b$ 的夹角为\mbox{（　　）}

（A）$\dfrac23\pi$；（B）$\dfrac\pi3$；（C）$\dfrac\pi6$；（D）$0$。

\Needspace{4\baselineskip}
\textbf{解析}

\Needspace{5\baselineskip}
设 $|\boldsymbol a|=t$，$|\boldsymbol b|=2t$，$\boldsymbol a$ 与 $\boldsymbol b$ 的夹角为 $\theta$，记 $c=\cos\theta$，则
\nopagebreak[4]
\[\boldsymbol a\cdot\boldsymbol a=t^2,\qquad \boldsymbol b\cdot\boldsymbol b=4t^2,\qquad \boldsymbol a\cdot\boldsymbol b=2t^2c .\]
\Needspace{5\baselineskip}
两组向量各由 $2$ 个 $\boldsymbol a$ 与 $2$ 个 $\boldsymbol b$ 排成。若 $\boldsymbol x$ 中有 $k$ 个 $\boldsymbol a$ 与 $\boldsymbol y$ 中同位置的 $\boldsymbol a$ 配对（$k=0,1,2$），则同时有 $k$ 个 $\boldsymbol b$ 与 $\boldsymbol b$ 配对，剩下 $4-2k$ 对是 $\boldsymbol a$ 与 $\boldsymbol b$ 的混合配对。于是
\nopagebreak[4]
\[S_k=kt^2+4kt^2+(4-2k)\cdot2t^2c
=t^2\left[5k+4c(2-k)\right].\]
\Needspace{5\baselineskip}
三类取值为
\nopagebreak[4]
\[S_0=8ct^2,\qquad S_1=(5+4c)t^2,\qquad S_2=10t^2 .\]
\Needspace{5\baselineskip}
因为 $c\leqslant1<\dfrac54$，所以
\nopagebreak[4]
\[8c\leqslant5+4c,\qquad 8c\leqslant10,\]
\Needspace{5\baselineskip}
最小值总是 $S_0=8ct^2$。由题设最小值为 $4t^2$：
\nopagebreak[4]
\[8c=4\ \Longrightarrow\ c=\frac12\ \Longrightarrow\ \theta=\frac\pi3 .\]
此时三个值为 $4t^2$，$7t^2$，$10t^2$，最小值确为 $4t^2$。

选 B。

\gkanswer{答案：B。三类配对的值依次为 $8|\boldsymbol a|^2\cos\theta$、$(5+4\cos\theta)|\boldsymbol a|^2$、$10|\boldsymbol a|^2$，最小者为 $8|\boldsymbol a|^2\cos\theta=4|\boldsymbol a|^2$，故 $\cos\theta=\dfrac12$。}

\end{answer}
\begin{answer}
\gksolhead{11}{}

\textbf{题目}：

$\left(\dfrac{16}{81}\right)^{-\frac34}+\log_3\dfrac54+\log_3\dfrac45=$\gkblank{}。

\Needspace{4\baselineskip}
\textbf{解析}

\Needspace{5\baselineskip}
先算分数指数幂：
\nopagebreak[4]
\[\left(\frac{16}{81}\right)^{-\frac34}=\left(\frac{81}{16}\right)^{\frac34}
=\left(\left(\frac32\right)^4\right)^{\frac34}=\left(\frac32\right)^3=\frac{27}8 .\]
\Needspace{5\baselineskip}
再算两个对数：同底对数相加等于真数之积的对数，
\nopagebreak[4]
\[\log_3\frac54+\log_3\frac45=\log_3\left(\frac54\times\frac45\right)=\log_31=0 .\]
\Needspace{5\baselineskip}
故原式
\nopagebreak[4]
\[\frac{27}8+0=\frac{27}8 .\]

\gkanswer{答案：$\dfrac{27}8$。$\left(\dfrac{16}{81}\right)^{-\frac34}=\left(\dfrac32\right)^3=\dfrac{27}8$，两个对数之和为 $\log_31=0$。}

\end{answer}
\begin{answer}
\gksolhead{12}{}

\textbf{题目}：

如图，在等腰直角三角形 $ABC$ 中，斜边 $BC=2\sqrt2$，过点 $A$ 作 $BC$ 的垂线，垂足为 $A_1$；过点 $A_1$ 作 $AC$ 的垂线，垂足为 $A_2$；过点 $A_2$ 作 $A_1C$ 的垂线，垂足为 $A_3$；$\cdots$，依此类推。设 $BA=a_1$，$AA_1=a_2$，$A_1A_2=a_3$，$\cdots$，$A_5A_6=a_7$，则 $a_7=$\gkblank{}。
\nopagebreak[4]
\begin{center}% ============================================================
% f2014_anhui_arts_q12_perpendiculars.tex
% 2014 安徽卷（文科）第 12 题：逐次作垂线
% 来源 PDF：普通高考/2014/2014安徽文.pdf
% 原页：第 1 页（A3 横排）中栏下部，印面 859，第 12 题题干下方
% 说明：等腰直角三角形 ABC，斜边 BC=2√2，A 为直角顶点。
%       A1 为 A 到 BC 的垂足；A2 为 A1 到 AC 的垂足；
%       A3 为 A2 到 A1C 的垂足；A4 为 A3 到 A2C 的垂足；
%       A5 为 A4 到 A3C 的垂足；A6 为 A5 到 A4C 的垂足。
% ============================================================
\begin{tikzpicture}[>=Stealth, line width=0.7pt, font=\small,
  x=1cm, y=1cm, line cap=round, line join=round]

  \coordinate (B) at (-4.243,0);
  \coordinate (C) at (4.243,0);
  \coordinate (A) at (0,4.243);
  \coordinate (A1) at (0,0);
  \coordinate (A2) at (2.121,2.121);
  \coordinate (A3) at (2.121,0);
  \coordinate (A4) at (3.182,1.061);
  \coordinate (A5) at (3.182,0);
  \coordinate (A6) at (3.712,0.530);

  % 三角形与逐次垂线
  \draw (B) -- (C) -- (A) -- cycle;
  \draw (A) -- (A1);
  \draw (A1) -- (A2);
  \draw (A2) -- (A3);
  \draw (A3) -- (A4);
  \draw (A4) -- (A5);
  \draw (A5) -- (A6);

  % 顶点标注
  \node[above] at (A) {$A$};
  \node[below left] at (B) {$B$};
  \node[below right] at (C) {$C$};
  \node[below] at (A1) {$A_1$};
  \node[right] at (A2) {$A_2$};
  \node[below] at (A3) {$A_3$};
  \node[right] at (A4) {$A_4$};
  \node[below] at (A5) {$A_5$};
  \node[right] at (A6) {$A_6$};

\end{tikzpicture}
\end{center}
\nopagebreak[4]
\Needspace{4\baselineskip}
\textbf{解析}

\Needspace{5\baselineskip}
等腰直角三角形的斜边 $BC=2\sqrt2$，故两直角边
\nopagebreak[4]
\[AB=AC=\frac{BC}{\sqrt2}=2 .\]
\Needspace{5\baselineskip}
于是
\nopagebreak[4]
\[a_1=BA=2,\qquad a_2=AA_1=\frac{BC}2=\sqrt2 .\]

\Needspace{5\baselineskip}
以下说明每作一次垂线，所得线段与上一条之比恒为 $\dfrac{\sqrt2}2$。以 $A_1A_2$ 为例：$A_1$ 是斜边中点，$AA_1\perp BC$，而 $A_2$ 是 $A_1$ 到 $AC$ 的垂足，故 $\triangle AA_1C$ 与 $\triangle A_1A_2C$ 都是等腰直角三角形，且后者与前者的相似比为 $\dfrac{A_1C}{AC}=\dfrac{\sqrt2}{2}$，于是
\nopagebreak[4]
\[a_3=A_1A_2=\frac{\sqrt2}2\cdot AA_1=\frac{\sqrt2}2\cdot\sqrt2=1 .\]
\Needspace{5\baselineskip}
同理，每次新作的垂线段都把当前的等腰直角三角形分成一个与之相似、相似比为 $\dfrac{\sqrt2}2$ 的小等腰直角三角形，故
\nopagebreak[4]
\[a_{n+1}=\frac{\sqrt2}2\,a_n\quad(n\geqslant1).\]
\Needspace{5\baselineskip}
所以 $\{a_n\}$ 是首项 $2$、公比 $\dfrac{\sqrt2}2$ 的等比数列：
\nopagebreak[4]
\[a_n=2\left(\frac{\sqrt2}2\right)^{n-1}.\]
\Needspace{5\baselineskip}
取 $n=7$：
\nopagebreak[4]
\[a_7=2\left(\frac{\sqrt2}2\right)^6=2\cdot\frac{1}{8}=\frac14 .\]

逐项核对：$a_3=1$，$a_4=\dfrac{\sqrt2}2$，$a_5=\dfrac12$，$a_6=\dfrac{\sqrt2}4$，$a_7=\dfrac14$，与上式一致。

\gkanswer{答案：$\dfrac14$。每一步所得的垂线段与上一条之比恒为 $\dfrac{\sqrt2}2$，$a_7=2\left(\dfrac{\sqrt2}2\right)^6=\dfrac14$。}

\end{answer}
\begin{answer}
\gksolhead{13}{}

\textbf{题目}：

不等式组 $\begin{cases}x+y-2\geqslant0,\\x+2y-4\leqslant0,\\x+3y-2\geqslant0\end{cases}$ 表示的平面区域的面积为\gkblank{}。

\Needspace{4\baselineskip}
\textbf{解析}

\Needspace{5\baselineskip}
三条边界直线为
\nopagebreak[4]
\[l_1:x+y-2=0,\qquad l_2:x+2y-4=0,\qquad l_3:x+3y-2=0 .\]
两两求交点：

\begin{itemize}
\item $l_1$ 与 $l_2$：相减得 $y=2$，代回得 $x=0$，交点为 $(0,2)$；
\item $l_1$ 与 $l_3$：相减得 $2y=0$，即 $y=0$，代回得 $x=2$，交点为 $(2,0)$；
\item $l_2$ 与 $l_3$：相减得 $y=-2$，代回得 $x=8$，交点为 $(8,-2)$。
\end{itemize}

\Needspace{5\baselineskip}
取可行域内的点 $(1,1)$ 检验：$1+1-2=0\geqslant0$，$1+2-4=-1\leqslant0$，$1+3-2=2\geqslant0$，故可行域是以这三点为顶点的三角形（含边界）。用行列式求面积：
\nopagebreak[4]
\[\begin{aligned}
S&=\frac12\left|x_1(y_2-y_3)+x_2(y_3-y_1)+x_3(y_1-y_2)\right|\\
&=\frac12\left|0\cdot(0+2)+2\cdot(-2-2)+8\cdot(2-0)\right|=\frac12\times8=4 .
\end{aligned}\]

\gkanswer{答案：$4$。三条边界两两交于 $(0,2)$，$(2,0)$，$(8,-2)$，用行列式计算得面积 $\dfrac12\times8=4$。}

\end{answer}
\begin{answer}
\gksolhead{14}{}

\textbf{题目}：

\Needspace{5\baselineskip}
若函数 $f(x)$（$x\in\mathbf R$）是周期为 $4$ 的奇函数，且在 $[0,2]$ 上的解析式为
\nopagebreak[4]
\[f(x)=\begin{cases}x(1-x),&0\leqslant x\leqslant1,\\\sin\pi x,&1<x\leqslant2,\end{cases}\]
则 $f\left(\dfrac{29}4\right)+f\left(\dfrac{41}6\right)=$\gkblank{}。

\Needspace{4\baselineskip}
\textbf{解析}

由周期为 $4$ 与 $f$ 为奇函数，先把自变量化到 $[0,2]$ 上。

\Needspace{5\baselineskip}
对 $f\left(\dfrac{29}4\right)$：因为 $\dfrac{29}4=7+\dfrac14$，连续减去两个周期 $8$ 得
\nopagebreak[4]
\[\frac{29}4-8=-\frac34,\]
\Needspace{5\baselineskip}
故
\nopagebreak[4]
\[f\left(\frac{29}4\right)=f\left(-\frac34\right)=-f\left(\frac34\right).\]
\Needspace{5\baselineskip}
而 $\dfrac34\in[0,1]$，用第一段解析式：
\nopagebreak[4]
\[f\left(\frac34\right)=\frac34\left(1-\frac34\right)=\frac34\cdot\frac14=\frac3{16},\]
\Needspace{5\baselineskip}
所以
\nopagebreak[4]
\[f\left(\frac{29}4\right)=-\frac3{16}.\]

\Needspace{5\baselineskip}
对 $f\left(\dfrac{41}6\right)$：$\dfrac{41}6=6+\dfrac56$，减去两个周期 $8$ 得
\nopagebreak[4]
\[\frac{41}6-8=-\frac76,\]
\Needspace{5\baselineskip}
故
\nopagebreak[4]
\[f\left(\frac{41}6\right)=f\left(-\frac76\right)=-f\left(\frac76\right).\]
\Needspace{5\baselineskip}
而 $\dfrac76\in(1,2]$，用第二段解析式：
\nopagebreak[4]
\[f\left(\frac76\right)=\sin\frac{7\pi}6=-\frac12,\]
\Needspace{5\baselineskip}
所以
\nopagebreak[4]
\[f\left(\frac{41}6\right)=\frac12 .\]
\Needspace{5\baselineskip}
相加得
\nopagebreak[4]
\[f\left(\frac{29}4\right)+f\left(\frac{41}6\right)=-\frac3{16}+\frac12=\frac5{16}.\]

\gkanswer{答案：$\dfrac5{16}$。$f\left(\dfrac{29}4\right)=f\left(-\dfrac34\right)=-\dfrac3{16}$，$f\left(\dfrac{41}6\right)=f\left(-\dfrac76\right)=-\sin\dfrac{7\pi}6=\dfrac12$。}

\end{answer}
\begin{answer}
\gksolhead{15}{}

\textbf{题目}：

若直线 $l$ 与曲线 $C$ 满足下列两个条件：（i）直线 $l$ 在点 $P(x_0,y_0)$ 处与曲线 $C$ 相切；（ii）曲线 $C$ 在点 $P$ 附近位于直线 $l$ 的两侧，则称直线 $l$ 在点 $P$ 处“切过”曲线 $C$。

下列命题正确的是\gkblank{}。（写出所有正确命题的编号）

① 直线 $l:y=0$ 在点 $P(0,0)$ 处“切过”曲线 $C:y=x^3$；

② 直线 $l:x=-1$ 在点 $P(-1,0)$ 处“切过”曲线 $C:y=(x+1)^2$；

③ 直线 $l:y=x$ 在点 $P(0,0)$ 处“切过”曲线 $C:y=\sin x$；

④ 直线 $l:y=x$ 在点 $P(0,0)$ 处“切过”曲线 $C:y=\tan x$；

⑤ 直线 $l:y=x-1$ 在点 $P(1,0)$ 处“切过”曲线 $C:y=\ln x$。

\Needspace{4\baselineskip}
\textbf{解析}

“切过”要求同时满足：$l$ 在 $P$ 处与 $C$ 相切，且 $C$ 在 $P$ 附近分居 $l$ 两侧。逐条判断。

① $C:y=x^3$ 在 $P(0,0)$ 处导数为 $0$，切线正是 $y=0$；且 $x<0$ 时 $x^3<0$，$x>0$ 时 $x^3>0$，曲线在 $l$ 两侧，符合。

② $C:y=(x+1)^2$ 在 $P(-1,0)$ 处导数为 $2(x+1)\big|_{x=-1}=0$，切线是水平线 $y=0$，不是 $x=-1$；$x=-1$ 甚至不是曲线的切线，不符合。

③ $C:y=\sin x$ 在 $P(0,0)$ 处导数为 $\cos0=1$，切线为 $y=x$。令 $h(x)=\sin x-x$，则 $h'(x)=\cos x-1\leqslant0$ 且仅 $x=0$ 时取等，$h$ 严格递减，$h(0)=0$，故 $x<0$ 时 $h>0$、$x>0$ 时 $h<0$，曲线在 $y=x$ 两侧，符合。

④ $C:y=\tan x$ 在 $P(0,0)$ 处导数为 $\sec^20=1$，切线为 $y=x$。在 $\left(-\dfrac\pi2,\dfrac\pi2\right)$ 内 $\tan x-x$ 的导数为 $\sec^2x-1=\tan^2x\geqslant0$，故 $\tan x-x\geqslant0$，仅 $x=0$ 取等，曲线只在 $y=x$ 的一侧，不符合。

⑤ $C:y=\ln x$ 在 $P(1,0)$ 处导数为 $\dfrac1x\bigg|_{x=1}=1$，切线为 $y=x-1$。令 $\varphi(x)=\ln x-(x-1)$，则 $\varphi'(x)=\dfrac1x-1$，$\varphi$ 在 $(0,1)$ 递增、在 $(1,+\infty)$ 递减，最大值为 $\varphi(1)=0$，故 $\ln x\leqslant x-1$，曲线只在切线一侧，不符合。

故正确的是 ①③。

\gkanswer{答案：①③。①中 $x^3$、③中 $\sin x-x$ 在 $P$ 两侧变号；②的 $x=-1$ 不是切线，④中 $\tan x-x$ 两侧同号，⑤中 $\ln x-(x-1)\leqslant0$ 只在单侧。}

\end{answer}
\begin{answer}
\gksolhead{16}{}

\textbf{题目}：

设 $\triangle ABC$ 的内角 $A$，$B$，$C$ 所对边的长分别是 $a$，$b$，$c$，且 $b=3$，$c=1$，$\triangle ABC$ 的面积为 $\sqrt2$，求 $\cos A$ 与 $a$ 的值。

\Needspace{4\baselineskip}
\textbf{解析}

\Needspace{5\baselineskip}
由面积公式
\nopagebreak[4]
\[S=\frac12bc\sin A=\frac12\times3\times1\times\sin A=\frac32\sin A=\sqrt2,\]
\Needspace{5\baselineskip}
得
\nopagebreak[4]
\[\sin A=\frac{2\sqrt2}3 .\]
\Needspace{5\baselineskip}
于是
\nopagebreak[4]
\[\cos A=\pm\sqrt{1-\sin^2A}=\pm\sqrt{1-\frac89}=\pm\frac13 .\]
两种符号都要检验是否构成三角形。

\Needspace{5\baselineskip}
情形一：$\cos A=\dfrac13$。由余弦定理
\nopagebreak[4]
\[a^2=b^2+c^2-2bc\cos A=9+1-2\times3\times1\times\frac13=8,\]
\Needspace{5\baselineskip}
故
\nopagebreak[4]
\[a=2\sqrt2 .\]
此时 $b=3<1+2\sqrt2$，三边可构成三角形。

\Needspace{5\baselineskip}
情形二：$\cos A=-\dfrac13$。同理
\nopagebreak[4]
\[a^2=9+1-2\times3\times1\times\left(-\frac13\right)=12,\]
\Needspace{5\baselineskip}
故
\nopagebreak[4]
\[a=2\sqrt3 .\]
此时 $b=3<1+2\sqrt3$，三边同样可构成三角形。

\Needspace{5\baselineskip}
所以
\nopagebreak[4]
\[\cos A=\frac13,\ a=2\sqrt2\quad\text{\text{或}}\quad \cos A=-\frac13,\ a=2\sqrt3 .\]

\gkanswer{答案：$\cos A=\dfrac13$，$a=2\sqrt2$；或 $\cos A=-\dfrac13$，$a=2\sqrt3$。由 $S=\dfrac12bc\sin A$ 得 $\sin A=\dfrac{2\sqrt2}3$，故 $\cos A=\pm\dfrac13$，再用余弦定理。}

\end{answer}
\begin{answer}
\gksolhead{17}{}

\textbf{题目}：

某高校共有学生 $15000$ 人，其中男生 $10500$ 人，女生 $4500$ 人，为调查该校学生每周平均体育运动时间的情况，采用分层抽样的方法，收集 $300$ 位学生每周平均体育运动时间的样本数据（单位：小时）。
\nopagebreak[4]
\begin{center}% ============================================================
% f2014_anhui_arts_q17_histogram.tex
% 2014 安徽卷（文科）第 17 题：频率分布直方图
% 来源 PDF：普通高考/2014/2014安徽文.pdf
% 原页：第 1 页（A3 横排）右栏，印面 859，第 17 题题干下方
% 说明：组距 2，六组 [0,2]，(2,4]，(4,6]，(6,8]，(8,10]，(10,12]，
%       频率/组距依次为 0.025，0.100，0.150，0.125，0.075，0.025，
%       各组频率之和为 0.05+0.20+0.30+0.25+0.15+0.05=1。
% ============================================================
\begin{tikzpicture}[>=Stealth, line width=0.7pt, font=\small,
  x=1.05cm, y=1cm, line cap=round, line join=round]

  % 纵坐标：数值放大 30 倍（0.150 -> 4.5cm）
  % 坐标轴
  \draw[->] (0,0) -- (7.20,0);
  \node[right, inner sep=2pt] at (6.42,0.30) {时间（小时）};
  \draw[->] (0,0) -- (0,5.45);
  \node[right, inner sep=1pt, align=center] at (0.10,4.95)
    {$\dfrac{\text{频率}}{\text{组距}}$};

  % 六个矩形
  \draw (0,0) rectangle (1.05,0.75);
  \draw (1.05,0) rectangle (2.10,3.00);
  \draw (2.10,0) rectangle (3.15,4.50);
  \draw (3.15,0) rectangle (4.20,3.75);
  \draw (4.20,0) rectangle (5.25,2.25);
  \draw (5.25,0) rectangle (6.30,0.75);

  % 虚线与纵轴刻度
  \foreach \yy/\lab in {0.75/0.025, 2.25/0.075, 3.00/0.100, 3.75/0.125, 4.50/0.150}{
    \draw[dashed, line width=0.4pt] (0,\yy) -- (6.90,\yy);
    \node[left, inner sep=3pt] at (0,\yy) {$\lab$};
  }
  \node[below, inner sep=5pt] at (0,0) {$0$};

  % 横轴刻度
  \foreach \xx/\lab in {1.05/2, 2.10/4, 3.15/6, 4.20/8, 5.25/10, 6.30/12}{
    \draw (\xx,0) -- (\xx,-0.10);
    \node[below, inner sep=5pt] at (\xx,0) {$\lab$};
  }

\end{tikzpicture}
\end{center}
\nopagebreak[4]
（1） 应收集多少位女生样本数据？

（2） 根据这 $300$ 个样本数据，得到学生每周平均体育运动时间的频率分布直方图（如图所示），其中样本数据分组区间为 $[0,2]$，$(2,4]$，$(4,6]$，$(6,8]$，$(8,10]$，$(10,12]$。估计该校学生每周平均体育运动时间超过 $4$ 个小时的概率；

（3） 在样本数据中，有 $60$ 位女生的每周平均体育运动时间超过 $4$ 个小时。请完成每周平均体育运动时间与性别的列联表，并判断是否有 $95\%$ 的把握认为“该校学生的每周平均体育运动时间与性别有关”。

\Needspace{10\baselineskip}
附：$K^2=\dfrac{n(ad-bc)^2}{(a+b)(c+d)(a+c)(b+d)}$

\begin{center}
\begin{tabular}{lllll}
\toprule
$P(K^2\geqslant k_0)$ & $0.10$ & $0.05$ & $0.010$ & $0.005$ \\
\midrule
$k_0$ & $2.706$ & $3.841$ & $6.635$ & $7.879$ \\
\bottomrule
\end{tabular}
\end{center}

\Needspace{4\baselineskip}
\textbf{解析}

\gksubpart{（1）}

\Needspace{5\baselineskip}
分层抽样按比例分配，女生在全校所占比例为
\nopagebreak[4]
\[\frac{4500}{15000}=\frac3{10},\]
\Needspace{5\baselineskip}
故女生样本量为
\nopagebreak[4]
\[300\times\frac3{10}=90 .\]

\gksubpart{（2）}

各组频率等于“频率/组距”乘以组距 $2$。由直方图读数：

\begin{center}
\begin{tabular}{lllllll}
\toprule
分组 & $[0,2]$ & $(2,4]$ & $(4,6]$ & $(6,8]$ & $(8,10]$ & $(10,12]$ \\
\midrule
频率/组距 & $0.025$ & $0.100$ & $0.150$ & $0.125$ & $0.075$ & $0.025$ \\
频率 & $0.05$ & $0.20$ & $0.30$ & $0.25$ & $0.15$ & $0.05$ \\
\bottomrule
\end{tabular}
\end{center}

频率之和 $0.05+0.20+0.30+0.25+0.15+0.05=1$，与直方图自洽。

\Needspace{5\baselineskip}
运动时间超过 $4$ 小时即落在后四组，频率为
\nopagebreak[4]
\[0.30+0.25+0.15+0.05=0.75 .\]
故估计概率为 $0.75$。

\gksubpart{（3）}

全校女生 $4500$ 人，样本中女生 $90$ 人，其中超过 $4$ 小时的有 $60$ 人，不超过的有 $30$ 人。样本中超过 $4$ 小时的总人数为 $300\times0.75=225$，其中男生 $225-60=165$ 人；男生样本共 $210$ 人，故不超过 $4$ 小时的男生为 $210-165=45$ 人。列联表为

\begin{center}
\begin{tabular}{llll}
\toprule
 & 超过 $4$ 小时 & 不超过 $4$ 小时 & 合计 \\
\midrule
男生 & $165$ & $45$ & $210$ \\
女生 & $60$ & $30$ & $90$ \\
合计 & $225$ & $75$ & $300$ \\
\bottomrule
\end{tabular}
\end{center}

\Needspace{5\baselineskip}
代入公式，$n=300$，$a=165$，$b=45$，$c=60$，$d=30$：
\nopagebreak[4]
\[K^2=\frac{300(165\times30-45\times60)^2}{210\times90\times225\times75}
=\frac{300\times2250^2}{318937500}
=\frac{100}{21}\approx4.762 .\]
\Needspace{5\baselineskip}
因为
\nopagebreak[4]
\[4.762>3.841,\]
所以在犯错误概率不超过 $0.05$ 的前提下，有 $95\%$ 的把握认为该校学生的每周平均体育运动时间与性别有关。

\gkanswer{答案（1）：$90$。按女生占比 $\dfrac{4500}{15000}=\dfrac3{10}$ 分层，$300\times\dfrac3{10}=90$。}

\gkanswer{答案（2）：$0.75$。后四组频率之和为 $2\times(0.150+0.125+0.075+0.025)=0.75$。}

\gkanswer{答案（3）：列联表为男生超过 $4$ 小时 $165$ 人、不超过 $45$ 人，女生分别为 $60$ 人与 $30$ 人；$K^2=\dfrac{100}{21}\approx4.762>3.841$，有 $95\%$ 的把握认为有关。}

\end{answer}
\begin{answer}
\gksolhead{18}{}

\textbf{题目}：

数列 $\{a_n\}$ 满足 $a_1=1$，$na_{n+1}=(n+1)a_n+n(n+1)$，$n\in\mathbf N^*$。

（1） 证明：数列 $\left\{\dfrac{a_n}n\right\}$ 是等差数列；

（2） 设 $b_n=3^n\sqrt{a_n}$，求数列 $\{b_n\}$ 的前 $n$ 项和 $S_n$。

\Needspace{4\baselineskip}
\textbf{解析}

\gksubpart{（1）}

\Needspace{5\baselineskip}
题设等号两边同除以 $n(n+1)$（$n\in\mathbf N^*$）：
\nopagebreak[4]
\[\frac{a_{n+1}}{n+1}=\frac{a_n}{n}+1,\]
\Needspace{5\baselineskip}
即
\nopagebreak[4]
\[\frac{a_{n+1}}{n+1}-\frac{a_n}{n}=1 .\]
\Needspace{5\baselineskip}
这说明数列 $\left\{\dfrac{a_n}n\right\}$ 从第 $1$ 项起，相邻两项之差恒为常数 $1$。又首项
\nopagebreak[4]
\[\frac{a_1}1=1,\]
故 $\left\{\dfrac{a_n}n\right\}$ 是以 $1$ 为首项、$1$ 为公差的等差数列。

\Needspace{5\baselineskip}
顺带写出通项：
\nopagebreak[4]
\[\frac{a_n}{n}=1+(n-1)\cdot1=n\ \Longrightarrow\ a_n=n^2 .\]

\gksubpart{（2）}

\Needspace{5\baselineskip}
由（1），$\sqrt{a_n}=\sqrt{n^2}=n$（$n\in\mathbf N^*$），故
\nopagebreak[4]
\[b_n=3^n\sqrt{a_n}=n\cdot3^n .\]
\Needspace{5\baselineskip}
于是
\nopagebreak[4]
\[S_n=1\cdot3+2\cdot3^2+3\cdot3^3+\cdots+n\cdot3^n .\]
\Needspace{5\baselineskip}
两边乘以 $3$：
\nopagebreak[4]
\[3S_n=1\cdot3^2+2\cdot3^3+\cdots+(n-1)\cdot3^n+n\cdot3^{n+1}.\]
\Needspace{5\baselineskip}
两式相减（错位相减）：
\nopagebreak[4]
\[S_n-3S_n=3+\left(3^2+3^3+\cdots+3^n\right)-n\cdot3^{n+1},\]
\Needspace{5\baselineskip}
即
\nopagebreak[4]
\[-2S_n=3+\frac{3^2\left(3^{n-1}-1\right)}{3-1}-n\cdot3^{n+1}
=3+\frac{3^{n+1}-9}{2}-n\cdot3^{n+1}.\]
\Needspace{5\baselineskip}
整理：
\nopagebreak[4]
\[-2S_n=\frac{6+3^{n+1}-9-2n\cdot3^{n+1}}2
=\frac{(1-2n)3^{n+1}-3}2,\]
\Needspace{5\baselineskip}
故
\nopagebreak[4]
\[S_n=\frac{(2n-1)3^{n+1}+3}4 .\]
检验 $n=1$：右边 $=\dfrac{1\cdot9+3}4=3=b_1$，正确。

\gkanswer{答案（1）：见解析版}

\gkanswer{答案（2）：$S_n=\dfrac{(2n-1)3^{n+1}+3}4$。由（1）得 $a_n=n^2$，$b_n=n\cdot3^n$，用错位相减法求和。}

\end{answer}
\begin{answer}
\gksolhead{19}{}

\textbf{题目}：

如图，四棱锥 $P-ABCD$ 的底面是边长为 $8$ 的正方形，四条侧棱长均为 $2\sqrt{17}$，点 $G$，$E$，$F$，$H$ 分别是棱 $PB$，$AB$，$CD$，$PC$ 上共面的四点，平面 $GEFH\perp$ 平面 $ABCD$，$BC\parallel$ 平面 $GEFH$。
\nopagebreak[4]
\begin{center}% ============================================================
% f2014_anhui_arts_q19_pyramid.tex
% 2014 安徽卷（文科）第 19 题：四棱锥 P-ABCD
% 来源 PDF：普通高考/2014/2014安徽文.pdf
% 原页：第 2 页（A3 横排）左栏下部，印面 860，第 19 题题干下方
% 说明：底面为边长 8 的正方形 ABCD，四条侧棱长 2√17，
%       顶点 P 在底面中心正上方，高为 6。
%       G，E，F，H 分别在 PB，AB，CD，PC 上，且 E 满足 EB=2，
%       平面 GEFH 垂直于底面，故其为过 EF（EF∥BC）的竖直平面。
%       被底面遮住的三条棱 PD、AD、DC 用虚线。
%       截面四边形 GEFH 中 GE、GH、HF 可见，EF 位于底面内被遮挡。
% ============================================================
\begin{tikzpicture}[>=Stealth, line width=0.7pt, font=\small,
  x=1cm, y=1cm, line cap=round, line join=round]

  % 底面斜二测画法：A(0,0), B(5.2,0), C(7.4,2.6), D(2.2,2.6)
  \coordinate (A) at (0,0);
  \coordinate (B) at (5.2,0);
  \coordinate (C) at (7.4,2.6);
  \coordinate (D) at (2.2,2.6);
  \coordinate (P) at (3.7,6.2);
  % E：AB 上 EB=2，AE=6
  \coordinate (E) at (3.9,0);
  % F：CD 上 CF=2
  \coordinate (F) at (6.1,2.6);
  % G，H 分别为 PB，PC 中点（平面为竖直平面 EF）
  \coordinate (G) at (4.45,3.1);
  \coordinate (H) at (5.55,4.4);

  % 不可见棱
  \draw[dashed] (P) -- (D);
  \draw[dashed] (A) -- (D);
  \draw[dashed] (D) -- (C);
  \draw[dashed] (E) -- (F);

  % 可见棱
  \draw (P) -- (A);
  \draw (P) -- (B);
  \draw (P) -- (C);
  \draw (A) -- (B);
  \draw (B) -- (C);

  % 截面四边形 GEFH（EF 已用虚线画出）
  \draw (G) -- (E);
  \draw (G) -- (H);
  \draw (H) -- (F);

  % 顶点标注
  \node[above] at (P) {$P$};
  \node[below left] at (A) {$A$};
  \node[below] at (B) {$B$};
  \node[right] at (C) {$C$};
  \node[above left, inner sep=2pt] at (D) {$D$};
  \node[below] at (E) {$E$};
  \node at (6.18,2.18) {$F$};
  \node at (4.02,3.10) {$G$};
  \node at (5.86,4.62) {$H$};

\end{tikzpicture}
\end{center}
\nopagebreak[4]
（1） 证明：$GH\parallel EF$；

（2） 若 $EB=2$，求四边形 $GEFH$ 的面积。

\Needspace{4\baselineskip}
\textbf{解析}

\gksubpart{（1）}

\Needspace{5\baselineskip}
因为 $BC\parallel$ 平面 $GEFH$，且 $BC\subset$ 平面 $PBC$，而平面 $PBC$ 与平面 $GEFH$ 的交线为 $GH$，由线面平行的性质定理得
\nopagebreak[4]
\[GH\parallel BC .\]
\Needspace{5\baselineskip}
又 $BC\subset$ 平面 $ABCD$，平面 $GEFH$ 与平面 $ABCD$ 的交线为 $EF$，且 $BC\parallel$ 平面 $GEFH$，同理得
\nopagebreak[4]
\[EF\parallel BC .\]
于是 $GH\parallel EF$。

\gksubpart{（2）}

\Needspace{5\baselineskip}
设底面中心为 $O$。正方形边长 $8$，故
\nopagebreak[4]
\[OA=OB=OC=OD=\frac{8\sqrt2}2=4\sqrt2 .\]
\Needspace{5\baselineskip}
四条侧棱相等，说明 $P$ 在底面的射影就是中心 $O$，设高 $PO=h$。在直角三角形 $POB$ 中
\nopagebreak[4]
\[PB^2=PO^2+OB^2\ \Longrightarrow\ (2\sqrt{17})^2=h^2+(4\sqrt2)^2,\]
\Needspace{5\baselineskip}
即
\nopagebreak[4]
\[68=h^2+32\ \Longrightarrow\ h=6 .\]

\Needspace{5\baselineskip}
建立坐标系：以 $A$ 为原点，$\overrightarrow{AB}$，$\overrightarrow{AD}$ 的方向分别为 $x$，$y$ 轴正方向，竖直向上为 $z$ 轴正方向，则
\nopagebreak[4]
\[A(0,0,0),\ B(8,0,0),\ C(8,8,0),\ D(0,8,0),\ P(4,4,6).\]
\Needspace{5\baselineskip}
由 $EB=2$ 得 $AE=6$，故 $E(6,0,0)$；又 $EF\parallel BC$ 且 $F$ 在 $CD$ 上，得 $F(6,8,0)$，于是
\nopagebreak[4]
\[EF=8 .\]

\Needspace{5\baselineskip}
平面 $GEFH$ 含底面上的直线 $EF$ 且垂直于底面，故它是过 $EF$ 的竖直平面，即 $x=6$。求它与 $PB$，$PC$ 的交点。直线 $PB$ 的参数式为
\nopagebreak[4]
\[(x,y,z)=(4,4,6)+t(4,-4,-6),\]
\Needspace{5\baselineskip}
令 $x=6$ 得 $t=\dfrac12$，故
\nopagebreak[4]
\[G=(6,2,3).\]
\Needspace{5\baselineskip}
同理直线 $PC$ 上令 $x=6$ 也得 $t=\dfrac12$，故
\nopagebreak[4]
\[H=(6,6,3).\]
\Needspace{5\baselineskip}
于是
\nopagebreak[4]
\[GH=4,\]
\Needspace{5\baselineskip}
且 $GH$ 与 $EF$ 都是水平线段，它们所在直线间的距离就是高度差
\nopagebreak[4]
\[d=3 .\]
\Needspace{5\baselineskip}
四边形 $GEFH$ 是上底 $GH=4$、下底 $EF=8$、高 $d=3$ 的梯形，面积
\nopagebreak[4]
\[S_{GEFH}=\frac{GH+EF}2\cdot d=\frac{4+8}2\times3=18 .\]

\gkanswer{答案（1）：见解析版}

\gkanswer{答案（2）：$18$。底面中心到各顶点距离 $4\sqrt2$，由侧棱 $2\sqrt{17}$ 得高 $6$；截面为梯形，$EF=8$，$GH=4$，高 $3$，面积 $\dfrac{8+4}2\times3=18$。}

\end{answer}
\begin{answer}
\gksolhead{20}{}

\textbf{题目}：

设函数 $f(x)=1+(1+a)x-x^2-x^3$，其中 $a>0$。

（1） 讨论 $f(x)$ 在其定义域上的单调性；

（2） 当 $x\in[0,1]$ 时，求 $f(x)$ 取得最大值和最小值时的 $x$ 的值。

\Needspace{4\baselineskip}
\textbf{解析}

\gksubpart{（1）}

\Needspace{5\baselineskip}
定义域为 $\mathbf R$，求导：
\nopagebreak[4]
\[f'(x)=(1+a)-2x-3x^2 .\]
\Needspace{5\baselineskip}
令 $f'(x)=0$，即 $3x^2+2x-(1+a)=0$，其判别式
\nopagebreak[4]
\[\Delta=4+12(1+a)=4(4+3a)>0,\]
\Needspace{5\baselineskip}
两根为
\nopagebreak[4]
\[x_1=\frac{-1-\sqrt{4+3a}}3,\qquad x_2=\frac{-1+\sqrt{4+3a}}3 .\]
\Needspace{5\baselineskip}
因为 $a>0$，有 $\sqrt{4+3a}>2$，故
\nopagebreak[4]
\[x_1<-1,\qquad x_2>\frac13 .\]
\Needspace{5\baselineskip}
$f'(x)=-3(x-x_1)(x-x_2)$ 是开口向下的二次函数，符号为：$x<x_1$ 时 $f'<0$，$x_1<x<x_2$ 时 $f'>0$，$x>x_2$ 时 $f'<0$。所以 $f(x)$ 在
\nopagebreak[4]
\[\left(-\infty,\frac{-1-\sqrt{4+3a}}3\right),\qquad
\left(\frac{-1+\sqrt{4+3a}}3,+\infty\right)\]
\Needspace{5\baselineskip}
上单调递减，在
\nopagebreak[4]
\[\left(\frac{-1-\sqrt{4+3a}}3,\frac{-1+\sqrt{4+3a}}3\right)\]
上单调递增。

\gksubpart{（2）}

\Needspace{5\baselineskip}
由（1），$f$ 在 $[0,1]$ 上的单调性取决于 $x_2$ 是否落在区间内：
\nopagebreak[4]
\[x_2=\frac{\sqrt{4+3a}-1}3 .\]
由 $x_2<1$ 得 $\sqrt{4+3a}<4$，即 $a<4$；且 $a>0$ 时 $x_2>\dfrac13>0$。

\Needspace{5\baselineskip}
\textbf{最大值。} 若 $0<a<4$，则 $x_2\in(0,1)$，$f$ 在 $[0,x_2]$ 增、在 $[x_2,1]$ 减，最大值在 $x=x_2$ 处取得：
\nopagebreak[4]
\[x=\frac{\sqrt{4+3a}-1}3 .\]
若 $a\geqslant4$，则 $x_2\geqslant1$，$f$ 在 $[0,1]$ 上单调递增，最大值在 $x=1$ 处取得。

\Needspace{5\baselineskip}
\textbf{最小值。} 只需比较端点。因为
\nopagebreak[4]
\[f(0)=1,\qquad f(1)=1+(1+a)-1-1=a,\]
所以

\begin{itemize}
\item 当 $0<a<1$ 时，$f(1)=a<1=f(0)$，最小值在 $x=1$ 处取得；
\item 当 $a=1$ 时，$f(0)=f(1)=1$，最小值在 $x=0$ 与 $x=1$ 处同时取得；
\item 当 $a>1$ 时，$f(0)=1<a=f(1)$，最小值在 $x=0$ 处取得。
\end{itemize}

\Needspace{5\baselineskip}
综上：最大值点为
\nopagebreak[4]
\[x=\frac{\sqrt{4+3a}-1}3\ (0<a<4),\qquad x=1\ (a\geqslant4);\]
\Needspace{5\baselineskip}
最小值点为
\nopagebreak[4]
\[x=1\ (0<a<1),\qquad x=0\ \text{\text{或}}\ x=1\ (a=1),\qquad x=0\ (a>1).\]

\gkanswer{答案（1）：记 $x_1=\dfrac{-1-\sqrt{4+3a}}3$，$x_2=\dfrac{-1+\sqrt{4+3a}}3$；减区间为 $(-\infty,x_1)$、$(x_2,+\infty)$，增区间为 $(x_1,x_2)$。}

\gkanswer{答案（2）：最大值点：$0<a<4$ 时为 $x=\dfrac{\sqrt{4+3a}-1}3$，$a\geqslant4$ 时为 $x=1$；最小值点：$0<a<1$ 时为 $x=1$，$a=1$ 时为 $x=0$ 或 $x=1$，$a>1$ 时为 $x=0$。}

\end{answer}
\begin{answer}
\gksolhead{21}{}

\textbf{题目}：

设 $F_1$，$F_2$ 分别是椭圆 $E:\dfrac{x^2}{a^2}+\dfrac{y^2}{b^2}=1$（$a>b>0$）的左、右焦点，过点 $F_1$ 的直线交椭圆 $E$ 于 $A$，$B$ 两点，$|AF_1|=3|BF_1|$。

（1） 若 $|AB|=4$，$\triangle ABF_2$ 的周长为 $16$，求 $|AF_2|$；

（2） 若 $\cos\angle AF_2B=\dfrac35$，求椭圆 $E$ 的离心率。

\Needspace{4\baselineskip}
\textbf{解析}

设焦距为 $2c$，则 $F_1(-c,0)$，$F_2(c,0)$，$c^2=a^2-b^2$。

\gksubpart{（1）}

\Needspace{5\baselineskip}
由 $|AF_1|=3|BF_1|$ 且 $A$，$F_1$，$B$ 共线、$F_1$ 在 $A$，$B$ 之间，得
\nopagebreak[4]
\[|AF_1|=\frac34|AB|=3,\qquad |BF_1|=\frac14|AB|=1 .\]
\Needspace{5\baselineskip}
由椭圆定义
\nopagebreak[4]
\[|AF_1|+|AF_2|=2a,\qquad |BF_1|+|BF_2|=2a,\]
\Needspace{5\baselineskip}
故
\nopagebreak[4]
\[|AF_2|=2a-3,\qquad |BF_2|=2a-1 .\]
\Needspace{5\baselineskip}
又 $\triangle ABF_2$ 的周长为 $16$：
\nopagebreak[4]
\[|AB|+|AF_2|+|BF_2|=4+(2a-3)+(2a-1)=16,\]
\Needspace{5\baselineskip}
即
\nopagebreak[4]
\[4a-4=12\ \Longrightarrow\ a=4 .\]
\Needspace{5\baselineskip}
于是
\nopagebreak[4]
\[|AF_2|=2a-3=5 .\]

\gksubpart{（2）}

\Needspace{5\baselineskip}
记 $|BF_1|=m$，则 $|AF_1|=3m$，$|AB|=4m$，且
\nopagebreak[4]
\[|AF_2|=2a-3m,\qquad |BF_2|=2a-m .\]
\Needspace{5\baselineskip}
设 $\angle AF_1F_2=\theta$，则 $\angle BF_1F_2=\pi-\theta$。在两个三角形中分别用余弦定理：
\nopagebreak[4]
\[|AF_2|^2=9m^2+4c^2-12mc\cos\theta, \tag{1}\]
\nopagebreak[4]
\[|BF_2|^2=m^2+4c^2+4mc\cos\theta . \tag{2}\]
\Needspace{5\baselineskip}
（1）减（2），左边用平方差：
\nopagebreak[4]
\[\left(2a-3m\right)^2-\left(2a-m\right)^2
=-2m\,(4a-4m)=-8m(a-m),\]
\Needspace{5\baselineskip}
右边为 $8m^2-16mc\cos\theta$。两边同除以 $8m>0$：
\nopagebreak[4]
\[-a+m=m-2c\cos\theta\ \Longrightarrow\ \cos\theta=\frac{a}{2c}\ \left(\geqslant\frac12\right). \tag{3}\]
\Needspace{5\baselineskip}
把（3）代入（1）：
\nopagebreak[4]
\[4a^2-12am+9m^2=9m^2+4c^2-6am
\ \Longrightarrow\ 6am=4(a^2-c^2),\]
\Needspace{5\baselineskip}
即
\nopagebreak[4]
\[m=\frac{2(a^2-c^2)}{3a} . \tag{4}\]
\Needspace{5\baselineskip}
于是
\nopagebreak[4]
\[|AF_2|=2a-3m=\frac{2c^2}a,\qquad
|BF_2|=2a-m=\frac{4a^2+2c^2}{3a},\qquad
|AB|=4m=\frac{8(a^2-c^2)}{3a} .\]
\Needspace{5\baselineskip}
记离心率 $e=\dfrac ca$，则 $c=ae$，三边化为
\nopagebreak[4]
\[|AF_2|=2ae^2,\qquad |BF_2|=\frac{a(4+2e^2)}3,\qquad |AB|=\frac{8a(1-e^2)}3 .\]
\Needspace{5\baselineskip}
在 $\triangle ABF_2$ 中，$\angle AF_2B$ 的对边是 $AB$，由余弦定理
\nopagebreak[4]
\[\cos\angle AF_2B
=\frac{|AF_2|^2+|BF_2|^2-|AB|^2}{2|AF_2|\cdot|BF_2|}=\frac35 .\]
\Needspace{5\baselineskip}
记 $s=e^2$，代入三边：
\nopagebreak[4]
\[|AF_2|^2+|BF_2|^2-|AB|^2
=a^2\left[4s^2+\frac{(4+2s)^2}9-\frac{64(1-s)^2}9\right]
=\frac{a^2}9\left(-24s^2+144s-48\right),\]
\nopagebreak[4]
\[2|AF_2|\cdot|BF_2|=2\cdot2as\cdot\frac{a(4+2s)}3=\frac{8a^2s(2+s)}3 .\]
\Needspace{5\baselineskip}
于是
\nopagebreak[4]
\[\frac{\dfrac{a^2}9\left(-24s^2+144s-48\right)}{\dfrac{8a^2s(2+s)}3}
=\frac{-s^2+6s-2}{s(2+s)}=\frac35 .\]
\Needspace{5\baselineskip}
去分母：
\nopagebreak[4]
\[5\left(-s^2+6s-2\right)=3s(2+s)
\ \Longrightarrow\ -5s^2+30s-10=6s+3s^2,\]
\Needspace{5\baselineskip}
即
\nopagebreak[4]
\[4s^2-12s+5=0\ \Longrightarrow\ s=\frac{12\pm\sqrt{144-80}}8=\frac{12\pm8}8,\]
\Needspace{5\baselineskip}
得
\nopagebreak[4]
\[s=\frac52\quad\text{\text{或}}\quad s=\frac12 .\]
\Needspace{5\baselineskip}
因为椭圆离心率 $e\in(0,1)$，故 $s=e^2<1$，只能取 $s=\dfrac12$，即
\nopagebreak[4]
\[e=\frac{\sqrt2}2 .\]

\Needspace{5\baselineskip}
验证：取 $a=\sqrt2$，$c=1$，$b=1$，则（3）给出 $\cos\theta=\dfrac{\sqrt2}2$，即 $\theta=\dfrac\pi4$；（4）给出 $m=\dfrac{2(2-1)}{3\sqrt2}=\dfrac{\sqrt2}3$。把 $F_1$ 平移到原点、$F_2$ 取 $(2,0)$，则 $A=\left(1,1\right)$，$B=\left(-\dfrac13,-\dfrac13\right)$，直接计算得 $|AF_1|=\sqrt2$，$|BF_1|=\dfrac{\sqrt2}3$，$|AF_2|=\sqrt2$，$|BF_2|=\dfrac{5\sqrt2}3$，两点均在椭圆 $\dfrac{(x-1)^2}2+y^2=1$ 上，且
\nopagebreak[4]
\[\cos\angle AF_2B=\frac35,\]
与题设相符。

\gkanswer{答案（1）：$5$。由 $|AB|=4$ 与 $|AF_1|=3|BF_1|$ 得 $|AF_1|=3$；周长给出 $|AF_2|+|BF_2|=12$，而两者分别为 $2a-3$ 与 $2a-1$，解得 $a=4$，$|AF_2|=5$。}

\gkanswer{答案（2）：$\dfrac{\sqrt2}2$。记 $|BF_1|=m$，由两个三角形余弦定理与 $|AF_1|=3m$ 得 $\cos\angle AF_1F_2=\dfrac a{2c}$、$m=\dfrac{2(a^2-c^2)}{3a}$，再对 $\triangle ABF_2$ 用余弦定理，得 $4e^4-12e^2+5=0$，取 $e<1$ 得 $e=\dfrac{\sqrt2}2$。}

\end{answer}